\section{Benchmarks And Evaluation Protocol}

The paper is written from a frozen official evidence surface with fixed datasets, model backbones, prompts, and scoring rules. The benchmark map below summarizes the benchmark set, and the cross-benchmark summary table that follows previews the empirical law that emerges from it.

\begin{table}[h]
\centering
\scriptsize
\begin{tabular}{p{1.8cm}p{1.8cm}p{1.9cm}p{1.0cm}p{1.4cm}p{3.1cm}}
\toprule
Benchmark & Regime & Source & Size & Models / seeds & Main question \\
\midrule
PopQA & Stable recurrence & Cached public PopQA slice & 500 & 1 / 5 & Can delayed utility replace repeated retrieval on recurring facts? \\
FreshQA main & Volatile public knowledge & Weekly public releases & 796 & 4 / 5 & Do changing answers favor retrieval-dominant behavior? \\
FreshQA sequence & Correction chains & Weekly public releases & 796 & 4 / 5 & Does the same boundary survive repeated corrections and forgetting probes? \\
MQuAKE & Propagation & Princeton-NLP MQuAKE-CF-3k & 21000 & 2 / 5 & Do edited facts propagate to single-hop and multi-hop consequences? \\
UniEdit & Broad editing & qizhou/UniEdit sampled test set & 3750 & 4 / 5 & Does the boundary survive a broad-domain editing benchmark? \\
\bottomrule
\end{tabular}

\caption{Benchmark map for the official benchmark surface. The headline evidence comes from frozen runs on PopQA, FreshQA public main and sequence, MQuAKE, and UniEdit.}
\label{tab:benchmark-map}
\end{table}

\begin{table}[h]
\centering
\scriptsize
\begin{tabular}{p{1.6cm}p{1.7cm}p{1.6cm}p{3.8cm}p{2.6cm}}
\toprule
Regime & Benchmark & Best policy & Key evidence & Law component \\
\midrule
Stable recurrence & PopQA & Router ties retrieval & 0.290 versus 0.290, retrieval $100.0 \rightarrow 72.2$, and no-$V_{\mathrm{future}}$ falls to 0.090 & Delayed utility helps when reuse is stable. \\
Volatile public knowledge & FreshQA main & Always retrieve & 1.000 on all 4 models; calibrated router 0.694--0.794; no-$V_{\mathrm{future}}$ 0.984--0.996 & Retrieval should dominate changing-answer public data. \\
Public correction chains & FreshQA sequence & Always retrieve & 1.000 on all 4 models; calibrated router 0.724--0.767; no-$V_{\mathrm{future}}$ 0.982--0.994 & Suppressing delayed utility is safer than soft calibration. \\
Propagation & MQuAKE & Always retrieve & 0.551 versus router 0.276--0.284; no-$V_{\mathrm{future}}$ 0.319--0.338 & Delayed utility does not transfer safely to propagation-heavy edits. \\
Broad editing sanity check & UniEdit & Near tie, retrieval best & Quality 0.998--1.000 across all models and modes & The benchmark is nondiagnostic for router superiority. \\
\bottomrule
\end{tabular}

\caption{Cross-benchmark summary from the frozen evidence surface. The key result is not broad router superiority; it is a boundary law showing where delayed utility helps and where retrieval should dominate.}
\label{tab:cross-benchmark-summary}
\end{table}

The benchmark surface has four regimes:
\begin{itemize}[leftmargin=1.5em]
\item \textbf{PopQA recurring-memory:} a cached recurring-heavy slice that asks whether reusable support can substitute for repeated retrieval on stable facts;
\item \textbf{FreshQA public main and sequence:} frozen public changing-answer evaluation derived from weekly releases, with the sequence slice stressing repeated corrections and forgetting probes;
\item \textbf{MQuAKE:} propagation-heavy knowledge editing where the important question is whether an updated fact transfers to single-hop and multi-hop consequences;
\item \textbf{UniEdit:} a broad editing benchmark used here as a high-coverage sanity check.
\end{itemize}

All runs use the same local practical model matrix:
\begin{itemize}[leftmargin=1.5em]
\item \texttt{google/flan-t5-small}
\item \texttt{google/flan-t5-base}
\item \texttt{Qwen/Qwen2.5-1.5B-Instruct}
\item \texttt{HuggingFaceTB/SmolLM2-360M-Instruct}
\end{itemize}

Main tables use five seeds \texttt{0,1,2,3,4} and only the allowed baselines:
\begin{itemize}[leftmargin=1.5em]
\item \texttt{always\_retrieve}
\item \texttt{retrieve\_gate} on FreshQA only
\item \texttt{router}
\item \texttt{router\_calibrated}
\item \texttt{router\_no\_v\_future}
\end{itemize}

This baseline set is deliberately controlled rather than leaderboard-oriented. \texttt{always\_retrieve} and \texttt{retrieve\_gate} anchor retrieval-dominant fallback behavior, \texttt{router} and \texttt{router\_calibrated} test the same six-action controller with and without public delayed-utility rescaling, and \texttt{router\_no\_v\_future} isolates the delayed-utility term while keeping the rest of the action inventory fixed. Across compared policies, we keep the backbone fixed within each table, expose the same timestamp-constrained evidence, use the same alias and stale-answer rules, and score every mode with the same evaluator. The support layer adds a 100-case manual FreshQA audit together with MQuAKE and UniEdit reliability notes, all cited later in the paper.

Some 95\% confidence intervals collapse to point intervals because all five official seeds produced identical aggregate scores for those settings, so the paired resampling layer had no observed variation to widen the interval.
